%% toc: P3 26 26.1,26.2,26.3
\section{Proofs: orders in Fra\"iss\'e's theory and source remarks of Fra\"iss\'e and Rado}
\label{sec:annex-p03-w03-p03-014-orders}

Orders are reflexive linear orders as in \cref{def:p03-w03-p03-014-eta}.
Throughout, an embedding of linear orders is the same as a strictly
increasing map: an injective map that preserves both \(\leq\) and its
negation is strictly increasing, and conversely a strictly increasing map between
linear orders is injective and satisfies \(x\leq y\iff f(x)\leq f(y)\),
since the orders are total.

\begin{annexitem}{annex:p03-w03-p03-014-chain-age}{thm:p03-w03-p03-014-chain-age}{The age of an infinite linear order}
\begin{proof}
Let \(R\) be an infinite linear order with base \(D\).

\emph{Members of the age are linear orders.} Reflexivity, antisymmetry,
transitivity and totality each concern at most three elements at a time
and are statements about the values of the relation on those elements.
They therefore hold in every restriction of \(R\) and in every
isomorphic copy of such a restriction. So each member of \(\gamma_R\)
is a finite linear order.

\emph{Every finite linear order is in the age.} Let \(A\) be a linear
order on a set of \(p\) elements. A finite linear order has a least
element, and removing it leaves a linear order, so by induction on \(p\)
its elements can be listed as \(a_1<\cdots<a_p\). As \(D\) is
infinite, choose \(p\) distinct elements of \(D\) and list them in the
same way as \(d_1<\cdots<d_p\). The bijection \(a_i\mapsto d_i\)
satisfies
\[
 A(a_i,a_j)=+\iff i\leq j\iff R(d_i,d_j)=+,
\]
so it is an isomorphism of \(A\) onto \(R\mathbin{\upharpoonright}\{d_1,\ldots,d_p\}\).
For \(p=0\) the empty relation is the restriction of \(R\) to the empty
set. Hence \(\gamma_R\) is exactly the class of finite linear orders.

\emph{The local class.} By \cref{def:p03-fraisse54-age},
\(A\in\Gamma_R\) means that every finite restriction of \(A\) embeds
into \(R\); since \(\gamma_R\) is closed under isomorphism, this says
that every finite restriction of \(A\) lies in \(\gamma_R\). If \(A\)
is a linear order, its finite restrictions are finite linear orders and
so lie in \(\gamma_R\). Conversely, if every finite restriction of
\(A\) is a linear order, then each instance of the four order axioms in
\(A\) involves at most three elements, holds in the restriction to those
elements, and therefore holds in \(A\). Thus \(\Gamma_R\) is the class of
linear orders.
\end{proof}
\end{annexitem}

\begin{annexitem}{annex:p03-w03-p03-014-ordering}{thm:p03-w03-p03-014-ordering}{Ordering an arbitrary set}
\begin{proof}
Let \(R\) be the order \((\mathbb N,\leq)\), which is infinite, and let
\(E\) be any set. By the size alternative in part~2 of
\cref{thm:p03-fraisse54-gamma}, either \(\Gamma_R\) consists of one
finite relation \(F\) together with its restrictions and their copies,
or \(\Gamma_R\) has a member whose base is \(E\).

The first alternative is impossible. Every relation in it has at most as
many elements as the base of \(F\), whereas by
\cref{thm:p03-w03-p03-014-chain-age} \(\Gamma_R\) contains the
linear orders of every finite size. So there is \(A\in\Gamma_R\) with
base \(E\), and by \cref{thm:p03-w03-p03-014-chain-age} again
\(A\) is a linear order on \(E\).

Choice enters only through the size alternative, whose proof takes a
normal ultrafilter on the maps \(E\to\mathbb N\) and its inverse
projection of \(R\) (\cref{thm:p03-fraisse54-local-projection}).
\end{proof}
\end{annexitem}

\begin{annexitem}{annex:p03-w03-p03-014-four-facts}{thm:p03-w03-p03-014-four-facts}{The four order propositions}
\begin{proof}
Fix an infinite order \(R\). By \cref{thm:p03-w03-p03-014-chain-age},
\(\Gamma_R\) is the class of linear orders.

\emph{Part~1.} Linear orders \(A,B\) lie in \(\Gamma_R\). By part~1 of
\cref{thm:p03-fraisse54-gamma}, \(\Gamma_R\) satisfies
\(\mathrm C_2\), so some \(C\in\Gamma_R\) has \(A\ll C\) and
\(B\ll C\); \(C\) is a linear order. For the explicit witness, replace
\(A,B\) by isomorphic copies \(A',B'\) with disjoint bases \(F,G\), and
form \(C=A'+B'\) on \(F\cup G\). It is reflexive because \(A'\) and
\(B'\) are. It is total: two elements on the same side are compared
there, and an element of \(F\) lies below one of \(G\). It is
antisymmetric: if \(x\leq y\leq x\), then \(x,y\) lie on the same side,
because no element of \(G\) lies below an element of \(F\), and that
side is antisymmetric. It is transitive: for \(x\leq y\leq z\), if all
three lie on one side, use that side; otherwise \(x\in F\) and
\(z\in G\), since an element of \(G\) is never below one of \(F\), and
then \(x\leq z\) by definition. The restriction of \(C\) to \(F\) is
\(A'\), so the inclusion of \(F\) composed with the copying isomorphism
is an embedding \(A\to C\); likewise for \(B\).

\emph{Part~2.} A set of linear orders is a set-indexed family of members
of \(\Gamma_R\). Part~3 of \cref{thm:p03-fraisse54-gamma} supplies a
member \(C\) of \(\Gamma_R\), hence a linear order, into which every
member of the family embeds.

\emph{Part~3.} The order \(\eta\) is nonempty, dense and without
endpoints, so it satisfies \(\mathcal H\) by
\cref{thm:p03-fraisse54-orders}. It is countable, hence homogeneous by
part~4 of \cref{thm:p03-fraisse54-extension}; in particular part~2 of
that theorem applies to it. Since \(\eta\) is infinite,
\(\Gamma_\eta\) is the class of all linear orders. Thus a finite or
countable linear order \(A\) satisfies \(A\prec\eta\), and part~2 of
\cref{thm:p03-fraisse54-extension} gives \(A\ll\eta\).

\emph{Part~4.} Let \(L\) be a nonempty countable dense order without
endpoints. A one-element order has its element as first and last element,
so \(L\) has at least two elements. It is infinite: starting from any
element, absence of a last element gives a strictly increasing sequence.
By \cref{thm:p03-fraisse54-orders}, \(L\) and \(\eta\) both satisfy
\(\mathcal H\). By \cref{thm:p03-w03-p03-014-chain-age}, both have the
finite linear orders as their age. Part~3 of
\cref{thm:p03-fraisse54-extension}, applied to the empty partial
isomorphism, yields an isomorphism \(L\to\eta\). This is the uniqueness
clause of \cref{thm:p03-fraisse54-original-construction} for the class
of linear orders, whose countable homogeneous representative is
therefore \(\eta\).

For the second sentence, let \(L\) be a countable homogeneous linear
order with more than one element. By part~1 of
\cref{thm:p03-fraisse54-extension} it satisfies \(\mathcal H\), and by
\cref{thm:p03-fraisse54-orders} it is dense without endpoints. By the
first sentence \(L\) is isomorphic to \(\eta\), and \(\eta\) is
homogeneous by the argument for part~3.
\end{proof}
\end{annexitem}

\begin{annexitem}{annex:p03-w03-p03-014-direct-gap}{thm:p03-w03-p03-014-direct-gap}{Homogeneous orders from the gap criterion}
\begin{proof}
The gap criterion of \cref{thm:p03-fraisse54-orders} says: a linear
order is homogeneous exactly when all its initial segments
\(\{x:x<a\}\) have one type, all its final segments \(\{x:x>a\}\) have
one type, and all its open intervals \(\{x:a<x<b\}\), \(a<b\), have
one type.

\emph{Part~1.} Let \(I\) be homogeneous with at least two elements.
Suppose \(b\) is a first element, and pick \(c\neq b\), so \(b<c\).
The initial segment below \(b\) is empty, and the one below \(c\)
contains \(b\). They have different types, contrary to the criterion.
The case of a last element is symmetric, using final segments. Now
suppose \(a<b\) with nothing between them. Since \(b\) is not last,
choose \(c>b\). The interval between \(a\) and \(b\) is empty, whereas
the interval between \(a\) and \(c\) contains \(b\), again contrary to
the criterion. So \(I\) is dense.

\emph{Part~2.} Let \(L\) be countable and dense, with at least two
elements and no endpoints. Let \(J\) be an initial segment below \(a\).
It is nonempty because \(a\) is not first. It has no first element: if
\(m\) were first in \(J\), an element below \(m\) in \(L\) would lie in
\(J\). It has no last element: if \(m\in J\), density gives an element
strictly between \(m\) and \(a\), which lies in \(J\). It is dense,
because an element of \(L\) between two members of \(J\) is again below
\(a\). It is countable, and a nonempty dense order without endpoints is
infinite. By \cref{thm:p03-rn-dlo-uniqueness}, which proves the countable
identification by alternating finite extensions and does not use
\(\mathcal H\), \(J\) is isomorphic to \(\eta\). The same argument, with
the order reversed, applies to final segments. For an interval between
\(a<b\): it is nonempty by density; an element \(m\) of it is neither
first nor last, since density gives elements between \(a\) and \(m\) and
between \(m\) and \(b\); and it is dense. So every interval is also
isomorphic to \(\eta\). All three families therefore have a single
type, and the gap criterion makes \(L\) homogeneous.
\end{proof}
\end{annexitem}

\begin{annexitem}{annex:p03-w03-p03-014-contains-eta}{thm:p03-w03-p03-014-contains-eta}{An infinite homogeneous order contains the rationals}
\begin{proof}
By part~1 of \cref{thm:p03-fraisse54-extension}, \(I\) satisfies
\(\mathcal H\). Since \(I\) is infinite,
\cref{thm:p03-w03-p03-014-chain-age} shows that \(\Gamma_I\) is the
class of all linear orders. A finite or countable linear order \(A\)
therefore satisfies \(A\prec I\), and part~2 of
\cref{thm:p03-fraisse54-extension}, whose target need not be countable,
gives \(A\ll I\). Taking \(A=\eta\) gives \(\eta\ll I\).
The argument used homogeneity only to obtain \(\mathcal H\).
\end{proof}
\end{annexitem}

\begin{annexitem}{annex:p03-w03-p03-014-condition-c}{thm:p03-w03-p03-014-condition-c}{Countable universal hosts for countable orders}
\begin{proof}
Suppose \(U\) is a countable linear order and \(\eta\ll U\). Let \(A\)
be a finite or countable linear order. By part~3 of
\cref{thm:p03-w03-p03-014-four-facts}, \(A\ll\eta\), and composing
with an embedding \(\eta\to U\) gives \(A\ll U\). So \(U\) witnesses
\(\mathcal C\). Taking \(U=\eta\) shows that the class of linear
orders satisfies \(\mathcal C\).

Conversely, let a finite or countable linear order \(U\) witness
\(\mathcal C\). Since \(\eta\) is a countable linear order,
\(\eta\ll U\). Then \(U\) contains an infinite subset, so it is
infinite, hence countably infinite.
\end{proof}
\end{annexitem}

\section{Proofs: additions to the Fra\"iss\'e presentation}
\label{sec:annex-p03-w03-p03-014-fraisse}

\begin{annexitem}{annex:p03-w03-p03-014-principal}{thm:p03-w03-p03-014-principal-filter}{Equality through a principal projection filter}
\begin{proof}
Let \(\varphi:E\to D\) and let \(\mathcal F_\varphi\) consist of all
sets \(Y\subseteq D^E\) with \(\varphi\in Y\). It contains \(D^E\),
is upward closed and closed under finite intersections, and omits the
empty set, so it is a proper filter. For \(a,b\in E\),
\[
 a\sim_{\mathcal F_\varphi}b
 \iff \{f\in D^E:f(a)=f(b)\}\in\mathcal F_\varphi
 \iff \varphi(a)=\varphi(b),
\]
the second equivalence because membership in \(\mathcal F_\varphi\)
means containing \(\varphi\).
\end{proof}
\end{annexitem}

\begin{annexitem}{annex:p03-w03-p03-014-printed-unary}{thm:p03-w03-p03-014-printed-unary}{The printed one-positive-point class and its repair}
\begin{proof}
For unary relations an embedding is an injection preserving the sign of
every element. Let \(\mathcal K_1\) be the class of unary relations with
exactly one positive element and \(\mathcal K_{\leq1}\) the class with
at most one.

\emph{\(\mathcal K_1\) fails \(\mathrm C_1\).} Take \(A\) on
\(\{a,c\}\) with \(A(a)=+\) and \(A(c)=-\). Then \(A\in\mathcal K_1\),
but its restriction to \(\{c\}\) has no positive element. Every class
\(\Gamma_R\) is closed under restriction
(\cref{thm:p03-fraisse54-gamma}, part~1), so \(\mathcal K_1\) is no
\(\Gamma_R\).

\emph{\(\mathcal K_1\) fails \(\overline{\mathcal D}\).} Let \(A\) on
\(\{a\}\) and \(B\) on \(\{b\}\), \(a\neq b\), both take the value
\(+\). Their bases are disjoint, so they agree on the common part. Any
relation extending both has the two positive elements \(a,b\), so no
common extension lies in \(\mathcal K_1\).

\emph{\(\mathcal K_1\) satisfies \(\mathcal D\) on finite members.}
Let \(f:C\to A\) and \(g:C\to B\) be embeddings with
\(A,B\in\mathcal K_1\) finite, with positive elements \(a\) and \(b\).
Let \(X\) be the set of points of the base of \(A\) outside
\(f[C]\cup\{a\}\). Take a set \(A'\) disjoint from the base of \(B\) and
a bijection \(x\mapsto x'\) from \(X\) onto it. Define \(M\) on the base
of \(B\) together with \(A'\) by \(M=B\) on the base of \(B\) and
\(M(x')=A(x)\) for \(x\in X\). Define \(e\) on the base of \(A\) by
\(e(f(c))=g(c)\) for \(c\) in the base of \(C\), by \(e(a)=b\) if
\(a\) is not in the image of \(f\), and by \(e(x)=x'\) for \(x\in X\).
If \(C\) has a positive element \(c_0\), then \(f(c_0)=a\) and
\(g(c_0)=b\), so the first clause already gives \(e(a)=b\). If not,
then \(a\notin f[C]\) and \(b\notin g[C]\), and the second clause sends
\(a\) to \(b\), which is not used by the first clause. In both cases
\(e\) is injective, \(e\circ f=g\), and \(e\) preserves signs: on
\(f[C]\) because \(A(f(c))=C(c)=B(g(c))\), at \(a\) because both values
are \(+\), and on the other points by the definition of \(M\). The
positive elements of \(M\) are \(b\) and the primed copies of positive
elements of \(X\), of which there are none because \(a\notin X\).
So \(M\in\mathcal K_1\), and the inclusion of \(B\) with \(e\) amalgamate
\(f\) and \(g\).

\emph{The repair.} Let \(R\) be a unary relation with exactly one
positive element and infinitely many negative elements. Its finite
restrictions are the finite unary relations with at most one positive
element, and every such relation occurs, so \(\gamma_R\) consists of the
finite members of \(\mathcal K_{\leq1}\); an arbitrary relation has at
most one positive element exactly when each of its finite restrictions
does, so \(\Gamma_R=\mathcal K_{\leq1}\). That its finite members
satisfy \(\mathcal D\) but not \(\overline{\mathcal D}\) is proved in
\cref{annex:p3.fraisse54.amalgamation}; the failure uses the same two
disjoint positive singletons as above.
\end{proof}
\end{annexitem}

\begin{annexitem}{annex:p03-w03-p03-014-arithmetic-types}{thm:p03-w03-p03-014-arithmetic-types}{Arithmetic types: countable hosts and a failure of condition C}
\begin{proof}
A relation \(A:E^n\to\{+,-\}\) is a structure for the one-symbol
language, and an injection between bases preserves both values of the
relation exactly when it preserves the atomic formulas and their
negations. So the embeddings of \cref{def:p03-fraisse54-relation} are the
embeddings of these structures. Let \(T\) be the set of sentences true in
\(A\); it is complete and satisfiable, and \(\mathcal K_A\) is the class
of models of \(T\).

\emph{Part~1.} \Cref{thm:p03-fraisse54-logical-joint} gives a model
\(N\) of \(T\) and embeddings \(f_i:M_i\to N\) of the members of the
family. Suppose the family is finite or countable and every \(M_i\) is
finite or countable. The union \(X\) of the images \(f_i[M_i]\) is then
finite or countable. By \cref{thm:downward-ls} there is an elementary
substructure \(N_0\preccurlyeq N\) containing \(X\) whose base has size
at most \(\max(\abs X,\aleph_0)\), the language being finite; so \(N_0\)
is finite or countable. It satisfies the same sentences as \(N\), so
\(N_0\in\mathcal K_A\). As a substructure, \(N_0\) carries the
restriction of the relation of \(N\), so each \(f_i\), with its values in
\(N_0\), is an embedding \(M_i\to N_0\). For the empty family, apply the
same theorem with \(X=\varnothing\).

\emph{The relation \(M\).} Let \(W\) be the countable set of finite
words over \(\{0,1\}\), the empty word included. The base of \(M\) is
\(W\times\mathbb Z\). For \((u,i)\) put \(b(u,i)=u_i\), the \(i\)-th
letter, when \(i\geq0\) is less than the length of \(u\), and \(b(u,i)=0\) otherwise. Let
\(M((u,i),(v,j))=+\) exactly when \(u=v\) and either \(j=i+1\), or
\(j=i\) and \(b(u,i)=1\). Write \(x\mathrel{R}y\) for the relation
symbol and \(T\) for the set of sentences true in \(M\).

The following two sentences are true in \(M\), hence belong to \(T\):
\[
\begin{gathered}
 \forall x\,\exists y\,\bigl(y\neq x\land x\mathrel Ry\land
   \forall z\,((z\neq x\land x\mathrel Rz)\to z=y)\bigr),\\
 \forall y\,\exists x\,\bigl(x\neq y\land x\mathrel Ry\land
   \forall z\,((z\neq y\land z\mathrel Ry)\to z=x)\bigr).
\end{gathered}
\]
Indeed the only element \(z\neq(u,i)\) with \((u,i)\mathrel Rz\) is
\((u,i+1)\), and the only \(z\neq(u,j)\) with \(z\mathrel R(u,j)\) is
\((u,j-1)\). Hence in every model \(U\) of \(T\) there are functions
\(s_U\) and \(p_U\): \(s_U(x)\) is the unique \(y\neq x\) with
\(x\mathrel Ry\), and \(p_U(y)\) the unique \(x\neq y\) with
\(x\mathrel Ry\). In \(M\), \(s_M(u,i)=(u,i+1)\) and
\(p_M(u,i)=(u,i-1)\). For \(d\) in \(U\) and \(k\in\mathbb Z\) put
\(d^{(k)}=s_U^k(d)\) for \(k\geq0\) and \(d^{(k)}=p_U^{-k}(d)\) for
\(k<0\), and define the mark word \(w_d:\mathbb Z\to\{0,1\}\) by
\(w_d(k)=1\) exactly when \(d^{(k)}\mathrel Rd^{(k)}\) holds in \(U\).

For \(k\geq0\) let \(\theta_k(x,y)\) be
\[
 \exists y_0\cdots\exists y_k\Bigl(y_0=x\land y_k=y\land
 \bigwedge_{j<k}\bigl(y_j\neq y_{j+1}\land y_j\mathrel Ry_{j+1}\bigr)\Bigr),
\]
and let \(\theta_{-k}(x,y)\) be the same formula with
\(y_{j+1}\mathrel Ry_j\) in place of \(y_j\mathrel Ry_{j+1}\). By
induction on \(k\), using the uniqueness clauses above, a model \(U\) of
\(T\) satisfies \(\theta_k(d,e)\) exactly when \(e=d^{(k)}\), for every
\(k\in\mathbb Z\).

\emph{One countable member for each mark word.} Fix
\(w:\mathbb Z\to\{0,1\}\). Add a constant \(c\) and let \(\Phi_w\)
consist of the sentences \(\exists y(\theta_k(c,y)\land y\mathrel Ry)\)
for \(w(k)=1\) and \(\exists y(\theta_k(c,y)\land\neg\,y\mathrel Ry)\)
for \(w(k)=0\), for all \(k\in\mathbb Z\). In a model of \(T\), these say
exactly that \(w_c=w\). A finite subset of \(T\cup\Phi_w\) mentions
only indices with \(-m\leq k\leq m\) for some \(m\geq0\). Let \(u\) be the word
\(w(-m)w(-m+1)\cdots w(m)\) of length \(2m+1\) and interpret \(c\) in
\(M\) as \((u,m)\). Then \(c^{(k)}=(u,m+k)\) and
\(b(u,m+k)=w(k)\) for \(-m\leq k\leq m\), so \(M\) satisfies the finite
subset. By compactness, \cref{thm:fol-compactness}, \(T\cup\Phi_w\) has a
model, and by \cref{thm:downward-ls} a countable one. Its reduct
\(N_w\) to the original language lies in \(\mathcal K_M\) and contains
an element \(c_w\) with \(w_{c_w}=w\). It is infinite because \(M\) is
and \(T\) contains, for each \(m\), a sentence saying that there are at
least \(m\) elements; so \(N_w\) is countably infinite.

\emph{Embeddings carry the mark word.} Let \(g:N\to U\) be an embedding
between members of \(\mathcal K_M\), and let \(d\in N\). For \(k\geq0\),
the elements \(d^{(k)}\neq d^{(k+1)}\) satisfy
\(d^{(k)}\mathrel Rd^{(k+1)}\) in \(N\). Since \(g\) is injective and
preserves \(R\), \(g(d^{(k)})\neq g(d^{(k+1)})\) and
\(g(d^{(k)})\mathrel Rg(d^{(k+1)})\) in \(U\), so
\(g(d^{(k+1)})=s_U(g(d^{(k)}))\). By induction
\(g(d^{(k)})=g(d)^{(k)}\) for \(k\geq0\); the same argument with
\(p\) gives it for \(k<0\). Since \(g\) preserves both values of the
relation, \(d^{(k)}\mathrel Rd^{(k)}\) holds in \(N\) exactly when
\(g(d)^{(k)}\mathrel Rg(d)^{(k)}\) holds in \(U\). Hence
\(w_{g(d)}=w_d\).

\emph{Failure of \(\mathcal C\).} Suppose \(U\in\mathcal K_M\) were
finite or countable and every finite or countable member of
\(\mathcal K_M\) embedded into it. For each
\(w:\mathbb Z\to\{0,1\}\) choose an embedding \(g_w:N_w\to U\). By the
previous paragraph \(w_{g_w(c_w)}=w\). So \(w\mapsto g_w(c_w)\) is
injective, because distinct words cannot be the mark word of one point.
There are \(2^{\aleph_0}\) such words, while \(U\) is finite or
countable, a contradiction. Thus \(\mathcal K_M\) fails \(\mathcal C\).
\end{proof}
\end{annexitem}

\section{Proofs: additions to the Rado presentation}
\label{sec:annex-p03-w03-p03-014-rado}

\begin{annexitem}{annex:p03-w03-p03-014-tarski-cf}{thm:p03-w03-p03-014-tarski-cf}{Tarski's cofinality index and ordinary cofinality}
\begin{proof}
Let \(\kappa=\aleph_n\), identified with the initial ordinal
\(\omega_n\). This is a nonzero limit ordinal, so by
\cref{thm:p06-cofinality-facts} its cofinality
\(\operatorname{cf}(\kappa)\) is an infinite regular cardinal, hence
\(\operatorname{cf}(\kappa)=\aleph_\nu\) for a unique ordinal \(\nu\).
Call a cardinal \(\mu\) \emph{sufficient} if \(\kappa\) is the sum of a
family of cardinals indexed by a set of cardinality \(\mu\), each
smaller than \(\kappa\). We show that \(\operatorname{cf}(\kappa)\) is
sufficient and that no smaller cardinal is; then the least ordinal
\(\nu'\) for which \(\aleph_{\nu'}\) is sufficient is \(\nu\), which is
the claim \(\operatorname{cf}(\aleph_n)=\aleph_{\mathrm{cf}(n)}\).

\emph{\(\operatorname{cf}(\kappa)\) is sufficient.} Choose a cofinal
subset of \(\kappa\) of order type \(\operatorname{cf}(\kappa)\) and
enumerate it as \((\alpha_i)_{i<\operatorname{cf}(\kappa)}\). Each
\(\alpha_i+1\) is below the limit ordinal \(\kappa\), so
\(\kappa_i=|\alpha_i+1|<\kappa\) because \(\kappa\) is an initial
ordinal. Every \(\beta<\kappa\) satisfies \(\beta\leq\alpha_i\) for some
\(i\), so \(\kappa=\bigcup_i(\alpha_i+1)\). With Choice,
\[
 \kappa\leq\sum_{i<\operatorname{cf}(\kappa)}\kappa_i
 \leq\operatorname{cf}(\kappa)\cdot\kappa=\kappa,
\]
and the sum equals \(\kappa\).

\emph{No smaller cardinal is sufficient.} Suppose
\(\kappa=\sum_{i\in I}\kappa_i\) with \(|I|=\mu<\operatorname{cf}(\kappa)\)
and every \(\kappa_i<\kappa\). If \(\kappa=\aleph_0\), then
\(\operatorname{cf}(\kappa)=\aleph_0\), so \(I\) is finite and all
\(\kappa_i\) are finite, and the sum is finite, a contradiction. If
\(\kappa>\aleph_0\), the set \(\{\kappa_i:i\in I\}\subseteq\kappa\) has
cardinality at most \(\mu\), so its order type is below the cardinal
\(\operatorname{cf}(\kappa)\) and it is not cofinal: some \(\beta<\kappa\)
exceeds every \(\kappa_i\). Put \(\lambda=\max(|\beta|,\mu,\aleph_0)\),
which is below \(\kappa\) since \(\mu<\operatorname{cf}(\kappa)\leq\kappa\)
and \(\kappa\) is uncountable. Then
\(\sum_{i\in I}\kappa_i\leq\mu\cdot|\beta|\leq\lambda<\kappa\), again
a contradiction.
\end{proof}
\end{annexitem}

\begin{annexitem}{annex:p03-w03-p03-014-rado-k2}{thm:p03-w03-p03-014-rado-k2}{Complexes as two-colour functions}
\begin{proof}
\emph{Graphs.} A graph \((A,B)\) has \(B\subseteq[A]^2\), which is the
definition of a one-dimensional complex in
\cref{def:p03-rado64-complex}. An injection \(\varphi\) is an embedding
of graphs exactly when \(\{x,y\}\in B\iff\{\varphi(x),\varphi(y)\}\in B^*\)
for \(x\neq y\), and this is the complex embedding condition for
\(2\)-element subsets. Two distinct \(2\)-element sets cannot share two
elements, so they meet in at most one point.

\emph{The dictionary.} Fix \(\ell\geq1\). For an \(\ell\)-dimensional
complex \(C=(\mathbb N_0,B)\) define \(\chi_C\in K_{2\ell}\) by
\(\chi_C(x_0,\ldots,x_\ell)=1\) if \(\{x_0,\ldots,x_\ell\}\in B\) and
\(0\) otherwise, for \(x_0<\cdots<x_\ell\). Since each
\((\ell+1)\)-element set has exactly one increasing enumeration, this is a
bijection between such complexes and \(K_{2\ell}\). For a strictly
increasing \(\varphi:\mathbb N_0\to\mathbb N_0\), the image of
\(\{x_0<\cdots<x_\ell\}\) is \(\{\varphi(x_0)<\cdots<\varphi(x_\ell)\}\),
so \(\varphi\) is an embedding \(C\to C'\) exactly when
\(\chi_C(\bar x)=\chi_{C'}(\varphi\bar x)\) for all increasing
\(\bar x\). The same construction with tuples of every length
\(\lambda+1\geq1\) is a bijection between mixed complexes on
\(\mathbb N_0\) and \(K_2\), with the same translation of embeddings.

\emph{The universal functions.} For \(k=2\) the function of
\cref{thm:p03-rado64-finite-colours} is
\(f^*_{2\ell}(x_0,\ldots,x_\ell)=d_2(x_\ell,\tau(x_0,\ldots,x_{\ell-1}))\),
which is \(\chi_{C^*_\ell}\) for the complex of
\cref{thm:p03-rado64-uniform-complex}. For \(\ell=1\),
\(\tau(x_0)=\binom{x_0}{1}=x_0\), so \(f^*_{21}(x_0,x_1)=d_2(x_1,x_0)\)
for \(x_0<x_1\), which is the adjacency of the graph \(R\) of
\cref{def:p03-rado64-digits}. Likewise \(f^*_2=\chi\) of the mixed complex
of \cref{thm:p03-rado64-mixed-complex}, as both are defined by the digit
\(d_2(x_\ell,\rho(x_0,\ldots,x_{\ell-1}))\).

\emph{Theorems 1, 3 and 4 from Theorem 5.} Let \(C\) be an
\(\ell\)-dimensional complex with at most countably many vertices.
Transport it along a bijection of its vertex set onto an initial segment
of \(\mathbb N_0\); if that segment is finite, add the remaining
integers as vertices lying in no simplex. Call the result \(C^\circ\).
Theorem 5 with \(k=2\) gives a strictly increasing \(\varphi\) with
\(\chi_{C^\circ}=f^*_{2\ell}\circ\varphi\); by the dictionary,
\(\varphi\) embeds \(C^\circ\) into \(C^*_\ell\) (into \(R\) when
\(\ell=1\)). Restricting \(\varphi\) to the transported original
vertices and composing with the transport embeds \(C\), because a
restriction of an embedding to a subset of vertices preserves and
reflects membership of the simplices on that subset, and the added
vertices lie in no simplex. The mixed case is the same with \(K_2\).
\end{proof}
\end{annexitem}
